% Source: https://artofproblemsolving.com/community/c3h1703016_vol_1_doubt
% Stable identity: aops-topic-1703016-q01
Use coordinates with side length $4$:
\[
B=(0,0),\quad C=(4,0),\quad A=(2,2\sqrt3).
\]
Then $M=\left(1,\sqrt3\right)$, $N=(1,0)$, $S=(2,0)$, $K=(3,0)$, and $P$, the midpoint of $CM$, is $\left(\dfrac{5}{2},\dfrac{\sqrt3}{2}\right)$.

\begin{asy}
usepackage("amsmath");
size(6.5cm);
defaultpen(fontsize(10pt));
real s=3^0.5;
pair B=(0,0),C=(4,0),A=(2,2*s);
pair M=(A+B)/2,Nn=(1,0),Sp=(2,0),K=(3,0),P=(M+C)/2;
draw(A--B--C--cycle,linewidth(0.9));
draw(C--M,linewidth(0.6));
draw(M--Nn,blue+linewidth(0.7));draw(P--Nn,blue+linewidth(0.7));draw(P--K,blue+linewidth(0.7));
real q=0.22;
draw(Nn+(0,q)--Nn+(q,q)--Nn+(q,0),linewidth(0.5));
pair u=unit(Nn-P),v=unit(K-P);
draw(P+q*u--P+q*(u+v)--P+q*v,linewidth(0.5));
dot(M);dot(Nn);dot(Sp);dot(K);dot(P);
label("$A$",A,N);label("$B$",B,SW);label("$C$",C,SE);
label("$M$",M,NW);label("$N$",Nn,S);label("$S$",Sp,S);label("$K$",K,S);label("$P$",P,NE);
label("\mbox{}",(min(currentpicture,user=true)+max(currentpicture,user=true))/2,black); // keeps the text after the figure black when LaTeX typesets the labels inline
\end{asy}

\textbf{$\angle MNB=90^\circ$.} $M$ and $N$ have the same $x$-coordinate, so $MN$ is vertical and perpendicular to $BC$. (Synthetically: $BM=2=2BN$ and $\angle B=60^\circ$, so triangle $MBN$ is a $30$-$60$-$90$ triangle with the right angle at $N$.)

\textbf{$\angle KPN=90^\circ$.} We have
\[
\overrightarrow{PK}=\left(\dfrac{1}{2},-\dfrac{\sqrt3}{2}\right),\qquad \overrightarrow{PN}=\left(-\dfrac{3}{2},-\dfrac{\sqrt3}{2}\right),
\]
and their dot product is $-\dfrac{3}{4}+\dfrac{3}{4}=0$. So $PK\perp PN$ and $\angle KPN=90^\circ$. (Synthetically: $NK=2$ and $PS=\sqrt{\dfrac{1}{4}+\dfrac{3}{4}}=1$, so $S$, the midpoint of $NK$, is at distance $1=\dfrac{NK}{2}$ from $P$, and $P$ lies on the circle with diameter $NK$.)

Hence $\angle MNB=\angle KPN=90^\circ$. $\blacksquare$
